\section{Introduction}
\label{sec:intro}

Metagenomic taxonomic classification---assigning each sequencing read to the
organism it originated from---is a foundational step in microbiome analysis,
underpinning studies of community composition, host--microbe interaction, and
clinical diagnostics \citep{Quince2017,HMP2012}. The dominant production tools
are alignment-free $k$-mer methods such as Kraken\,2 and Centrifuge, which match
long exact $k$-mers ($k\approx 31$--$35$) against a reference database and are
fast, accurate, and interpretable when the target organism is well represented
in that database \citep{Wood2019,Kim2016}. Their weaknesses are equally
well established: performance degrades sharply on organisms absent from or
poorly covered by the reference, and exact-match indexing cannot exploit the
distributional structure of sequence that a learned model can
\citep{Ye2019,Breitwieser2019,McIntyre2017}.

In parallel, genomic foundation models (GFMs)---large self-supervised
transformers pre-trained on genomic sequence---have advanced rapidly, with
Nucleotide Transformer, DNABERT/DNABERT-2, HyenaDNA, Caduceus, and Evo
demonstrating strong transfer to regulatory, variant-effect, and
genome-design tasks \citep{DallaTorre2023,Ji2021,Zhou2024dnabert2,
Nguyen2023hyena,Schiff2024caduceus,Nguyen2024,Brixi2025}. This success invites
an obvious question for the microbiome community: \emph{can, and should, GFMs be
used for metagenomic short-read taxonomic classification?} The setting is
adversarial for a foundation model. Reads are short ($\sim$150\,bp), the
taxonomic signal per read is faint, the label space is large and
long-tailed, and most GFMs were pre-trained on eukaryotic (often
human) genomes rather than the bacterial and archaeal diversity that
dominates real communities. Prior neural approaches such as DeepMicrobes,
BERTax, and MetaTransformer show that learned read classifiers are feasible
\citep{Liang2020,Mock2022bertax,Wichmann2023}, but they do not isolate
\emph{why} they work, and they predate the current generation of large
pre-trained backbones. More recently, ICCTax adapted the HyenaDNA genomic
language model to hierarchical taxonomic classification with an explicit focus on
out-of-distribution generalization \citep{Gao2025icctax}. Our focus is
complementary rather than competing: ICCTax classifies kilobase-scale contigs
(1{,}500\,bp) with a single purpose-built classifier, whereas we evaluate
150\,bp shotgun reads, benchmark several transformer GFMs rather than proposing
one model, and isolate whether pre-training, tokenization, or data scale
controls performance.

The practical decision facing a method developer is therefore not ``does a
GFM work at all'' but ``which design choice actually matters''. A GFM-based
classifier is defined by at least four largely independent axes: (i) how much
data it is trained on; (ii) whether the backbone is pre-trained or trained
from scratch; (iii) how the sequence is tokenized (the $k$-mer length and
whether tokens overlap); and (iv) the level at which the output is
evaluated---per read, or aggregated to a sample-level community profile.
These axes are usually entangled in published comparisons, so a reader cannot
tell whether a reported gain comes from a bigger model, more data, or a
different tokenizer. A recent study of tokenizer choice in genomic language
models \citep{Lindsey2025tokenizer} underscores that this axis alone can
dominate, but its effect has not been quantified against pre-training and data
scale in the metagenomic read-classification setting.

This paper answers the methodological question head-on with a controlled
benchmark that changes \emph{one factor at a time}. Holding architecture,
training data, and evaluation protocol fixed, we separately measure the effect
of data scale, pre-training, and tokenization, and we evaluate both at the
read level and at the sample level against Kraken\,2 (raw and with Bracken
re-estimation), with a preliminary real mock-community check. Our central finding is
that the axis most often left unexamined---tokenization---is the one that sets
the performance ceiling, while data scaling helps only until that ceiling is
reached and pre-training contributes a smaller but real gain once tokenization
is fixed. We then translate these findings into concrete, use-case-specific
recommendations for practitioners and for the design of the next generation of
microbial foundation models. Sections are organized as a problem-solving
protocol: we first explain why the task is hard
(Section~\ref{sec:challenges}), describe a leakage-controlled benchmark and a
dual read/sample evaluation (Section~\ref{sec:design}), then present the data
scaling (Section~\ref{sec:scaling}), pre-training-versus-tokenization
(Section~\ref{sec:decomp}), and read-versus-sample (Section~\ref{sec:sample})
results, before giving practical recommendations
(Section~\ref{sec:recommendations}) and limitations
(Section~\ref{sec:limitations}).
